%% ----------------------------------------------------------------------------
% CVG SA/MA thesis template
%
% Created 03/08/2024 by Tobias Fischer
%% ----------------------------------------------------------------------------

\vspace{3cm}

\chapter*{Abstract}
\noindent
Robots, vehicles, and wearable devices that repeatedly traverse the same environment, a fleet of cars driving through the same city, or a person wearing an AR headset through the same building produce their own trajectory estimate from onboard visual or visual-inertial odometry. These per-session trajectories are locally consistent but drift over time and live in unrelated, session-specific coordinate frames. This drift becomes increasingly worse with the scale of the area traversed. Turning many such independent, drifting trajectories into one globally consistent set of trajectories, where different sessions revisit the same place, is the large-scale multi-session trajectory optimization problem this thesis addresses.

\noindent
Existing solutions to this problem largely operate offline: every session's data must be collected and available up front, one session is often designated in advance as a fixed reference frame, and only then is a single fused pose graph solved for. This assumption breaks down for a continuously operating deployment, where sessions arrive one at a time, and there is no way to know in advance which one will end up anchoring the rest. In comparison to traditional SLAM methods, our approach operates entirely without a map: we optimize only the trajectories and absolute poses produced by learned pose estimation methods.

\noindent
This thesis presents Polygraph, a multi-session trajectory optimization pipeline that processes sessions online, one keyframe, and one spatial cluster at a time, without assuming that all sessions are available up front or that any one session is in a common reference frame. The odometry keyframes for each session are grouped into local spatial clusters. This odometry may come from any source, an agent's own on-device odometry, or any SLAM method agent might be running. When a cluster closes, its keyframes are matched via learned image retrieval against every other session's keyframes, and candidate cross-session pairs are verified through sparse feature matching, monocular metric depth lifting, and robust pose estimation. A two-tier consensus gate admits a cross-session loop-closure edge into the shared pose graph only once independent evidence corroborates it. This protects against a single spurious match corrupting the whole graph, while the graph itself is kept up to date through a cheap incremental solve after every cluster and periodically refined through a robust, outlier-pruning batch solve. In the end, even though input odometry can be visual-inertial, our method is purely visual.

\noindent
We evaluated Polygraph in multiple different setups chosen to be as different from each other as possible: stereo-camera outdoor driving (KITTI \cite{geiger2012kitti}), HoloLens visual-inertial indoor traversal (LaMAR \cite{sarlin2022lamar}), iPad visual-inertial indoor traversal (LaMAR \cite{sarlin2022lamar}), and Boston Dynamics Spot visual-inertial indoor legged locomotion (CroCoDL \cite{blum2025crocodl}), among others. In the one hardware-matched, head-to-head comparison against the closest prior work we could run (2GO's own slowest reported case, identical RTX 4090), Polygraph runs about three and a half times faster (586.57 s vs.\ 2067.31 s, a real-time factor of 37.9\% against their 133.5\%). In that same sequence, our method commits two orders of magnitude fewer loop-closure edges than the previous work \cite{lim2025twogo} (34 vs.\ 3403) while achieving comparable or better accuracy. We test multi-device multi-session capabilities with further experiments and show algorithm robustness to different viewpoints, internal algorithm parameters, and different camera setups. Together, these results show that robust multi-session trajectory optimization is achievable in a genuinely online setting, without dense 3D reconstruction, across a
substantially wider range of real-world conditions than previously demonstrated.